\documentclass[spanish,letterpaper,oneside]{article}

% Datos del documento
\def\documenttitle {¿Qué sabemos y cómo lo enseñamos?}
\def\documentsubtitle {Temas selectos de matemáticas I}
\def\documentsubject {Sesión 8 · Propuesta didáctica}
\def\documentauthor {Mtro. Martín Jonathan de la Cruz Muñoz}
\def\coursename {Temas selectos de matemáticas I}
\def\coursecode {PFM02}

% Identidad institucional
\def\universityname {Instituto de Investigación, Innovación y Estudios de Posgrado para la Educación del Estado de Nuevo León}
\def\universityfaculty {Maestría en Aprendizaje y Enseñanza de las Matemáticas}
\def\universitydepartment {Temas selectos de matemáticas I}
\def\universitydepartmentimage {img/iiiepe/marca-iiiepe.png}
\def\universitydepartmentimagecfg {height=1.75cm}
\def\universitylocation {Monterrey, Nuevo León}
\def\studentfullname {Mtro. Martín Jonathan de la Cruz Muñoz}
\def\studentid {00874}
\def\studentemail {martin.delacruz@campus.iiiepe.edu.mx}
\def\programperiod {1er tetramestre, marzo--julio de 2026}
\def\programmodality {Modalidad presencial}
\def\programshort {MAEM}

\def\authortable {
  \begin{tabular}{ll}
    \textbf{Alumno:} & \studentfullname \\
    Matrícula: & \studentid \\
    Correo institucional: & \studentemail \\
    Unidad: & \coursename \\
    Programa: & \programshort \\
    Periodo: & \programperiod \\
    & \\
    \multicolumn{2}{l}{Fecha de realización: 10 de octubre de 2026} \\
    \multicolumn{2}{l}{\universitylocation}
  \end{tabular}
}

\input{template}

\usepackage{helvet}
\renewcommand{\familydefault}{\sfdefault}
\usepackage{array}
\usepackage{longtable}
\usepackage{pdflscape}
\usepackage{tikz}
\usepackage[most]{tcolorbox}
\usetikzlibrary{arrows.meta,positioning,calc,fit,backgrounds}

\definecolor{iiiepeNavy}{HTML}{163B5C}
\definecolor{iiiepeBlue}{HTML}{1F6F9F}
\definecolor{iiiepeAqua}{HTML}{20B7A8}
\definecolor{iiiepeOrange}{HTML}{F28C28}
\definecolor{iiiepePaper}{HTML}{F6F8F7}
\definecolor{iiiepeInk}{HTML}{1E2A32}
\def\portraittitlecolor {iiiepeNavy}
\setcitestyle{authoryear,open={(},close={)}}

\newcommand{\insertcoverwatermark}{%
  \AddToShipoutPictureBG*{%
    \begin{tikzpicture}[remember picture,overlay]
      \node[opacity=.09,inner sep=0pt] at (current page.center)
      {\includegraphics[width=.58\paperwidth]{img/iiiepe/logo-puntos.png}};
    \end{tikzpicture}%
  }%
}

\newtcolorbox{consignabox}[1][]{
  enhanced,breakable,
  colback=iiiepePaper,colframe=iiiepeNavy,
  boxrule=1pt,leftrule=4pt,arc=2pt,
  left=8pt,right=8pt,top=6pt,bottom=6pt,
  before skip=6pt,after skip=10pt,
  fonttitle=\bfseries\color{white},coltitle=white,
  colbacktitle=iiiepeNavy,
  attach boxed title to top left={xshift=8pt,yshift=-2pt},
  boxed title style={colback=iiiepeNavy,colframe=iiiepeNavy,arc=1pt},
  title={Consigna de la actividad},#1}

\newtcolorbox{propuestabox}[1][]{
  enhanced,breakable,
  colback=iiiepePaper,colframe=iiiepeAqua,
  boxrule=1pt,leftrule=4pt,arc=2pt,
  left=8pt,right=8pt,top=6pt,bottom=6pt,
  before skip=6pt,after skip=10pt,
  fonttitle=\bfseries\color{white},coltitle=white,
  colbacktitle=iiiepeAqua,
  attach boxed title to top left={xshift=8pt,yshift=-2pt},
  boxed title style={colback=iiiepeAqua,colframe=iiiepeAqua,arc=1pt},
  title={Propuesta didáctica},#1}

\newtcolorbox{rubricabox}[1][]{
  enhanced,breakable,
  colback=iiiepePaper,colframe=iiiepeOrange,
  boxrule=1pt,leftrule=4pt,arc=2pt,
  left=8pt,right=8pt,top=6pt,bottom=6pt,
  before skip=6pt,after skip=10pt,
  fonttitle=\bfseries\color{white},coltitle=white,
  colbacktitle=iiiepeOrange,
  attach boxed title to top left={xshift=8pt,yshift=-2pt},
  boxed title style={colback=iiiepeOrange,colframe=iiiepeOrange,arc=1pt},
  title={Rúbrica de evaluación},#1}

\begin{document}

\insertcoverwatermark

% ============================================================
% PORTADA
% ============================================================
\begin{titlepage}
  \begin{center}
    \vspace*{1cm}
    
    \includegraphics[height=2.5cm]{img/iiiepe/marca-iiiepe.png}
    
    \vspace{0.5cm}
    
    {\Large\bfseries \universityname\par}
    
    \vspace{0.3cm}
    
    {\large \universityfaculty\par}
    
    \vspace{2cm}
    
    {\Huge\bfseries\color{iiiepeNavy} \documenttitle\par}
    
    \vspace{0.5cm}
    
    {\Large\color{iiiepeBlue} \documentsubject\par}
    
    \vspace{1cm}
    
    {\large \coursename\par}
    
    \vspace{0.3cm}
    
    {\large Bloque II: Enseñar desde el desafío\par}
    
    \vspace{2cm}
    
    \authortable
    
    \vfill
    
    {\large Sábado 10 de octubre de 2026\par}
    
  \end{center}
\end{titlepage}

% ============================================================
% CONSIGNA
% ============================================================
\section{Consigna de la actividad}

\begin{consignabox}
\textbf{Actividad 1. ¿Qué sabemos y cómo lo enseñamos?}

\vspace{0.3cm}

\textbf{Instrucciones presenciales:}
\begin{enumerate}
  \item Participa en el diagnóstico inicial para recuperar tus conocimientos previos sobre los contenidos que abordaremos.
  \item Resuelve los desafíos matemáticos utilizando los materiales disponibles. Explora distintas posibilidades, registra tus procedimientos y justifica tus conclusiones.
  \item Reflexiona sobre cada reto desde tu práctica docente. Identifica posibles dificultades de aprendizaje y plantea preguntas que ayuden a tus estudiantes a construir sus propias respuestas.
\end{enumerate}

\textbf{Instrucciones en línea:}

Diseña una propuesta didáctica en CANVA para enseñar el \textbf{área y el perímetro} en el grado escolar de tu elección. Considera:
\begin{itemize}
  \item Una situación desafiante
  \item Las estrategias y materiales que utilizarías
  \item Tu intervención como docente
  \item El aprendizaje que deseas promover
\end{itemize}

Comparte tu propuesta con el grupo y participa en el intercambio de ideas, reconociendo que un mismo contenido matemático puede enseñarse de distintas maneras.
\end{consignabox}

% ============================================================
% PROPUESTA DIDÁCTICA
% ============================================================
\section{Propuesta didáctica: Enseñanza del área y perímetro}

\begin{propuestabox}
\subsection{Datos generales}

\begin{itemize}
  \item \textbf{Grado escolar:} 5° de primaria (10--11 años)
  \item \textbf{Asignatura:} Matemáticas
  \item \textbf{Tema:} Área y perímetro de figuras geométricas
  \item \textbf{Duración estimada:} 3 sesiones de 50 minutos
  \item \textbf{Modalidad:} Presencial con apoyo de materiales concretos
\end{itemize}

\subsection{Situación desafiante: ``El huerto escolar''}

\textbf{Contexto:} La escuela ha decidido construir un huerto escolar en el patio. Los estudiantes de 5° grado han sido encargados de diseñar la distribución del espacio.

\textbf{El reto:}
\begin{quote}
``El director nos ha asignado un terreno rectangular de 12 metros de largo por 8 metros de ancho para construir nuestro huerto. Necesitamos diseñar la distribución considerando que queremos cultivar tomates, lechugas y zanahorias. Cada tipo de cultivo necesita un espacio diferente. Además, necesitamos cercar el huerto para protegerlo. ¿Cuánto material necesitaremos para la cerca? ¿Cómo distribuimos el espacio para cada cultivo?''
\end{quote}

\subsection{Estrategias didácticas}

\begin{enumerate}
  \item \textbf{Aprendizaje basado en proyectos:} Los estudiantes trabajarán en equipos para diseñar su propuesta de huerto escolar.
  \item \textbf{Uso de materiales concretos:} Geoplanos, cuerdas, reglas, papel cuadriculado y calculadoras.
  \item \textbf{Modelación matemática:} Representarán el terreno a escala y calcularán áreas y perímetros.
  \item \textbf{Discusión grupal:} Presentarán sus propuestas y justificarán sus decisiones.
\end{enumerate}

\subsection{Materiales}

\begin{itemize}
  \item Papel cuadriculado o cartulina
  \item Reglas y escalímetros
  \item Cuerdas o mecates
  \item Geoplanos
  \item Calculadoras
  \item Computadora o tablet para elaborar la presentación en Canva
\end{itemize}

\subsection{Intervención docente}

\textbf{Preguntas orientadoras (sin dar respuestas directas):}

\begin{enumerate}
  \item ``¿Qué información nos dan sobre el terreno? ¿Qué forma tiene?''
  \item ``Si queremos cercar el huerto, ¿qué medida necesitamos calcular? ¿Cómo la calculamos?''
  \item ``Si queremos saber cuánto espacio hay para cultivar, ¿qué medida necesitamos?''
  \item ``¿Es lo mismo el espacio para cultivar que la distancia alrededor del terreno?''
  \item ``¿Cómo podemos distribuir el espacio de manera equitativa para los tres cultivos?''
  \item ``¿Qué estrategias podemos usar para verificar nuestros cálculos?''
\end{enumerate}

\textbf{Anticipación de dificultades:}
\begin{itemize}
  \item Confundir área con perímetro
  \item No considerar las unidades de medida
  \item Dificultad para visualizar el problema a escala
  \item Errores en las operaciones con números decimales
\end{itemize}

\subsection{Aprendizaje esperado}

\textbf{Al finalizar la actividad, los estudiantes:}

\begin{itemize}
  \item Distinguirán entre área y perímetro como medidas diferentes de una figura geométrica.
  \item Calcularán el área y perímetro de rectángulos usando fórmulas matemáticas.
  \item Aplicarán estos conceptos en situaciones contextualizadas y significativas.
  \item Justificarán sus procedimientos y conclusiones usando el lenguaje matemático.
  \item Trabajarán colaborativamente para resolver problemas.
\end{itemize}

\subsection{Evidencia de aprendizaje}

\begin{itemize}
  \item Diseño del huerto escolar con cálculos de área y perímetro
  \item Presentación en Canva explicando su propuesta
  \item Exposición oral justificando las decisiones tomadas
\end{itemize}

\subsection{Enlace a la propuesta en Canva}

La propuesta didáctica visualizada en Canva está disponible en:
\begin{center}
  \url{https://www.canva.com/design/DAHXj5yH7Xg/QFhMML5K4wgghuVPSgpcaw/view}
\end{center}

\end{propuestabox}

% ============================================================
% RÚBRICA
% ============================================================
\section{Rúbrica de evaluación}

\begin{rubricabox}

\begin{longtable}{|p{3cm}|p{2cm}|p{2cm}|p{2cm}|p{2cm}|}
\hline
\textbf{Criterio} & \textbf{Excelente (10)} & \textbf{Bueno (9)} & \textbf{Aceptable (8)} & \textbf{Necesita mejorar (7)} \\
\hline
\endfirsthead

\hline
\textbf{Criterio} & \textbf{Excelente (10)} & \textbf{Bueno (9)} & \textbf{Aceptable (8)} & \textbf{Necesita mejorar (7)} \\
\hline
\endhead

1. Participación en los retos matemáticos & Participa activamente, explora distintas estrategias y justifica sus procedimientos y conclusiones. & Participa de manera constante, utiliza estrategias pertinentes y explica sus resultados. & Participa de forma limitada y resuelve los retos con explicaciones poco desarrolladas. & Muestra escasa participación y presenta dificultades para desarrollar o explicar sus procedimientos. \\
\hline
2. Reflexión pedagógica & Analiza críticamente los retos, anticipa dificultades de aprendizaje y propone intervenciones que favorecen el razonamiento autónomo. & Identifica dificultades de aprendizaje y plantea intervenciones docentes pertinentes. & Reconoce algunas dificultades, pero propone intervenciones generales o poco claras. & Presenta reflexiones superficiales, sin identificar dificultades ni estrategias de acompañamiento. \\
\hline
3. Grado escolar de la propuesta & Define claramente el grado escolar y adapta la actividad a las características y conocimientos de sus estudiantes. & Define el grado escolar y plantea una actividad generalmente adecuada a sus estudiantes. & Señala el grado escolar, pero la actividad muestra poca adaptación a sus características. & No precisa el grado escolar o plantea una actividad inadecuada para sus estudiantes. \\
\hline
4. Situación o reto matemático & Diseña una situación desafiante, contextualizada y clara que promueve la exploración y el razonamiento sobre área y perímetro. & Presenta un reto claro y pertinente que favorece el uso de conocimientos sobre área y perímetro. & Propone una situación relacionada con el contenido, pero centrada principalmente en aplicar procedimientos. & Plantea una actividad poco clara, descontextualizada o basada únicamente en ejercicios mecánicos. \\
\hline
5. Estrategia didáctica y materiales & Propone una estrategia creativa, viable y bien estructurada, con materiales que favorecen la participación y la construcción del aprendizaje. & Presenta una estrategia pertinente y materiales adecuados para desarrollar la actividad. & Describe una estrategia general, con materiales o procedimientos poco definidos. & No establece una estrategia clara o los materiales no contribuyen al propósito de aprendizaje. \\
\hline
6. Intervención docente & Anticipa dificultades específicas y formula preguntas orientadoras que promueven la autonomía, el análisis y la argumentación sin dar respuestas directas. & Reconoce posibles dificultades y propone preguntas o acciones pertinentes para orientar a los estudiantes. & Menciona dificultades generales, pero sus intervenciones ofrecen poca orientación para el razonamiento. & No anticipa dificultades o plantea intervenciones centradas en proporcionar directamente las respuestas. \\
\hline
7. Aprendizaje esperado & Establece con claridad qué comprensión matemática busca desarrollar y muestra coherencia entre el aprendizaje, el reto y la estrategia. & Define un aprendizaje pertinente y relacionado con la actividad propuesta. & Expresa un aprendizaje general o poco preciso, con una relación limitada con la actividad. & No define claramente el aprendizaje esperado o este no corresponde con la propuesta. \\
\hline
8. Entrega a tiempo & Entrega dentro del límite asignado posterior a la sesión. & Entrega un día posterior al límite asignado. & Entrega de 2 a 4 días posteriores al límite asignado. & Entrega de 5 días o más posteriores al límite asignado. \\
\hline

\end{longtable}

\end{rubricabox}

% ============================================================
% RECURSOS
% ============================================================
\section{Recursos de la plataforma}

\begin{itemize}
  \item \textbf{Consigna de la actividad:} Actividad 1\_Nombre\_Apellido.pdf
  \item \textbf{Rúbrica de evaluación:} Rúbrica Actividad 1.pdf
  \item \textbf{Presentación Canva:} \url{https://www.canva.com/design/DAHXj5yH7Xg/QFhMML5K4wgghuVPSgpcaw/view}
  \item \textbf{Plataforma:} \url{https://evaposgrados.iiiepe.edu.mx/aula/course/view.php?id=1049}
\end{itemize}

% ============================================================
% REFERENCIAS
% ============================================================
\section{Referencias}

\begin{itemize}
  \item Instituto de Investigación, Innovación y Estudios de Posgrado para la Educación del Estado de Nuevo León. (2026). \textit{Temas Selectos de Matemáticas I: Guía de actividades}. IIIEPE.
  \item SEP. (2023). \textit{Plan de Estudio para la Educación Primaria: Matemáticas}. Secretaría de Educación Pública.
  \item Van de Walle, J. A. (2018). \textit{Elementary and Middle School Mathematics: Teaching Developmentally}. Pearson.
\end{itemize}

\end{document}
